\section{Experiments}
\label{sec:experiments}

\subsection{Datasets and Baselines}

We evaluate on eight tasks from the SCROLLS benchmark~\cite{shaham2022scrolls}:
\textsc{QuALITY}, \textsc{QASPER}, \textsc{NarrativeQA}, \textsc{SQuALITY},
\textsc{Qasper-Abstract}, \textsc{GovReport}, \textsc{SumSCROLLS}, and
\textsc{ContractNLI}.
Baselines include Longformer-Base (4k tokens), BigBird-RoBERTa, LED, and PRIMERA.
All models are fine-tuned for 3 epochs on 4$\times$A100 GPUs with a batch size of 16.

\subsection{Main Results}

Table~\ref{tab:main} reports F\textsubscript{1} / ROUGE-L on the SCROLLS dev
split. \modelname{}-Base consistently improves over Longformer-Base by
\textbf{3.1 avg.\ F\textsubscript{1}} while reducing FLOPs per forward pass
by 18\%.

\begin{table}[!t]
\centering
\caption{\textbf{SCROLLS dev results.}
         F\textsubscript{1} (QA tasks) and ROUGE-L (summarisation tasks).
         $\Delta$ is the gap vs.\ Longformer-Base.
         \badge{Best} highlights top scores.}
\label{tab:main}
\setlength{\tabcolsep}{4pt}
\renewcommand{\arraystretch}{1.15}
\begin{tabularx}{\linewidth}{@{}l*{4}{>{\centering\arraybackslash}X}@{}}
\toprule
\rowcolor{headbg!8}
\textbf{Model} & \textbf{QuALITY} & \textbf{QASPER} & \textbf{GovRpt.} & \textbf{Avg.} \\
\midrule
Longformer-B    & 56.3 & 41.8 & 29.1 & 42.4 \\
BigBird-R       & 57.1 & 42.4 & 29.5 & 43.0 \\
LED-Base        & 55.9 & 41.1 & 30.2 & 42.4 \\
PRIMERA         & 58.6 & 43.7 & 31.8 & 44.7 \\
\midrule
\rowcolor{accentlt}
\modelname{}-B  & \textbf{60.1} & \textbf{45.6} & \textbf{33.2}
                & \textbf{45.5} \\
\quad $\Delta$  & \textcolor{accent}{+3.8} & \textcolor{accent}{+3.8}
                & \textcolor{accent}{+4.1} & \textcolor{accent}{+3.1} \\
\bottomrule
\end{tabularx}
\end{table}

\subsection{Ablation Studies}

Table~\ref{tab:ablation} isolates the contribution of each component. Removing
the context write ($k=0$) degrades to the Longformer baseline. Removing the
read gate collapses the bank to a fixed summary vector, losing 1.8 F\textsubscript{1}.
Increasing $k$ beyond 64 yields diminishing returns at higher FLOP cost.

\begin{table}[!t]
\centering
\caption{\textbf{Ablation on QuALITY dev.}
         All variants start from Longformer-Base.}
\label{tab:ablation}
\setlength{\tabcolsep}{4pt}
\renewcommand{\arraystretch}{1.15}
\begin{tabularx}{\linewidth}{@{}>{\raggedright\arraybackslash}Xl
                              >{\centering\arraybackslash}X
                              >{\centering\arraybackslash}X@{}}
\toprule
\rowcolor{headbg!8}
\textbf{Configuration} & & \textbf{F\textsubscript{1}} & \textbf{FLOPs (T)} \\
\midrule
No gateway ($k=0$)      & & 56.3 & 1.00 \\
Static bank (no read gate)& & 58.3 & 0.90 \\
\modelname{} $k=16$     & & 59.4 & 0.84 \\
\rowcolor{accentlt}
\modelname{} $k=32$     & & \textbf{60.1} & \textbf{0.82} \\
\modelname{} $k=64$     & & 60.3 & 0.97 \\
\bottomrule
\end{tabularx}
\end{table}

\subsection{Qualitative Analysis}

Figure~\ref{fig:attn} shows attention heat-maps on a \textsc{NarrativeQA}
example. The Context Gateway routes question-relevant passages to the bank
(dark cells in the write map), allowing all tokens to attend to them in the
read step---despite those passages being over 2\,000 tokens away.

% Attention heatmap placeholder figure
\begin{figure}[!t]
\centering
\begin{tikzpicture}
  % Background
  \fill[pagebg!50] (0,0) rectangle (7.2, 1.9);
  \draw[rule, line width=0.6pt] (0,0) rectangle (7.2, 1.9);

  % Left panel header
  \node[font=\tiny\bfseries\color{accent}] at (1.8, 1.75) {Write (top-$k$ scores)};
  % Right panel header
  \node[font=\tiny\bfseries\color{accent}] at (5.4, 1.75) {Read attention};

  % Divider
  \draw[rule, line width=0.4pt] (3.6, 0.1) -- (3.6, 1.65);

  % Left heatmap (schematic rectangles of varying darkness)
  \foreach \x/\y/\c in {
    0.2/1.4/30, 0.7/1.4/60, 1.2/1.4/90, 1.7/1.4/20, 2.2/1.4/10,
    2.7/1.4/80, 3.2/1.4/15,
    0.2/1.1/10, 0.7/1.1/20, 1.2/1.1/70, 1.7/1.1/40, 2.2/1.1/90,
    2.7/1.1/60, 3.2/1.1/30,
    0.2/0.8/50, 0.7/0.8/10, 1.2/0.8/20, 1.7/0.8/80, 2.2/0.8/30,
    2.7/0.8/15, 3.2/0.8/90,
    0.2/0.5/20, 0.7/0.5/90, 1.2/0.5/40, 1.7/0.5/10, 2.2/0.5/70,
    2.7/0.5/80, 3.2/0.5/50}{
    \fill[accent!\c!pagebg] (\x, \y) rectangle ++(0.38, 0.22);
  }

  % Right heatmap
  \foreach \x/\y/\c in {
    3.8/1.4/10, 4.3/1.4/30, 4.8/1.4/80, 5.3/1.4/90, 5.8/1.4/10,
    6.3/1.4/40, 6.8/1.4/20,
    3.8/1.1/20, 4.3/1.1/10, 4.8/1.1/90, 5.3/1.1/80, 5.8/1.1/30,
    6.3/1.1/10, 6.8/1.1/60,
    3.8/0.8/40, 4.3/0.8/20, 4.8/0.8/30, 5.3/0.8/60, 5.8/0.8/80,
    6.3/0.8/90, 6.8/0.8/10,
    3.8/0.5/10, 4.3/0.5/50, 4.8/0.5/20, 5.3/0.5/70, 5.8/0.5/90,
    6.3/0.5/30, 6.8/0.5/80}{
    \fill[accent!\c!pagebg] (\x, \y) rectangle ++(0.38, 0.22);
  }

  % Y-axis label
  \node[font=\tiny\color{neutralmid}, rotate=90] at (-0.05, 0.95) {Query token};
  % X-axis label (left)
  \node[font=\tiny\color{neutralmid}] at (1.8, 0.28) {Sequence position};
  \node[font=\tiny\color{neutralmid}] at (5.4, 0.28) {Context bank slot};

\end{tikzpicture}
\caption{\textbf{Attention visualisation.} Left: per-token write scores (darker
         = higher selection probability). Right: read attention over the $k=32$
         context bank. High-scoring writes cluster around question-relevant spans.}
\label{fig:attn}
\end{figure}

% ============================================================
